% !TEX root = ../math600-notes.tex

\section{Modules}

\begin{defn}
    Let $A$ be a ring. A (left) $A$-\textbf{module} is a set $(M,+,\cdot)$ endowed with two operations:
    $$+:M \times M \to M\:\: (x,y) \mapsto x+y$$
    $$\cdot :A \times M \to M \: \: (a,m) \mapsto a\cdot m$$
    such that for all $a,b \in A,x,y \in M$,
    \begin{enumerate}
        \item[(i)] $(M,+)$ is an additive abelian group
        \item[(ii)] a) $a(x+y)=ax+ay$, b) $(a+b)x = ax+bx$, c) $1\cdot x=x$, d) $(ab)\cdot x=a(bx)$
    \end{enumerate}
\end{defn}

\begin{rmk}
    Modules are to rings as vector spaces are to fields. Their significance is that modules "linearize" rings, and represent them. One can show that if $A$ is commutative, then the category $A$-MOD of all $A$ modules determines $A$.
\end{rmk}

\begin{thm}[Morita's Theorem]
    The category of $A$ modules is equivalent to the category of $M_n(A)$ modules, for all fixed $n \ge1$.
\end{thm}

\begin{ex}
    \begin{enumerate}
        \item $A$ is an $A$-module, in an obvious way. More generally, if $f:A \to B$ is any ring homomorphism, then $B$ becomes an $A$-module with product $a \cdot b:=f(a) \cdot b$ for $a\in A,b \in B$. In fact, we can replace $B$ by any $B$-module $N$, then $N$ becomes an $A$-module via the map $f$.
        \item If $A =F$ for $F$ a field, then an $A$-module is just an $F$-vector space.
        \item If $A = \Z$, then a $\Z$-module $M$ is the same as an abelian group, with $n\cdot x = (1+\overset{n}{\cdots}+1)x=x+\overset{n}{\cdots}+x$, $n \geq 1$.
        \item An example we'll study carefully later: Let $A = F[x]$, for $F$ a field, and $V$ be an $F[x]$-module. Then, $V$ is also an $F$-module, so $V$ is an $F$-vector space. The element $x$ induces an $F$-linear map $T:V \to V$ mapping $v \mapsto x\cdot v$. Now, any power $x^k$ acts on $V$ by $v \mapsto T^k(v)$. So, an $F[x]$-module is an $F$-vector space, together with a liner map $T:V \to V$. Conversely, given any $F$-vector space $V$, and any $F$-linear map $T:V \to V$, we can turn $V$ into an $F[x]$ module by setting $f(x) \cdot v= f(T)(v)$, for any $f \in F[x]$. Classifying the $F[x]$ modules up to isomorphism leads to the canonical forms of matrices.
    \end{enumerate}
\end{ex}

\textbf{From here on, we assume $A$ is commutative}

\begin{defn}
    \begin{enumerate}
        \item A subgroup $N \subset M$ is a \textbf{submodule} of $M$ if $N$ is an $A$-module with the induced operations from $M$.
        \item A map $\varphi: M \to N$ is a \textbf{homomorphism} of $A$-modules, or an $A$-linear map, if $\varphi(x+y)=\varphi(x)+\varphi(y)$, $\varphi(ax)=a\varphi(x)$, for all $x,y \in M, a \in A$.
        \item We define, for $\varphi :M \to N$ an $A$-linear map, $\ker \varphi$ and $\Img \varphi$ as for vector spaces. The first and second isomorphism theorems hold for modules.
        \item Given $N \subset M$ any submodule, we can always form the quotient module $M / N$ with the projection map $\pi : M \to M/N$. Submodules of $M/N$ correspond to submodules $M'$ of $M$ with $M' \subset N$
        \item Given any two $A$-modules $M$ and $N$, we have an $A$-module $\Hom _A(M,N)$ consisting of all $A$-linear maps $M \to N$. Taking $N=A$ defines the \textbf{dual module}.
        \item We can form sums and intersections of submodules of of an $A$-module $M$. If $S \subset M$ is any subset, we have the submodule $\angles{S}$ generated by $S$:
        $$\angles{S}:= \bigcap_{S \subset N\text{ submodule}}N= \{a_1s_1+\cdots + a_rs_r \mid a_i\in A,s_i \in S, r \geq 1\}.$$
        $M$ is \textbf{finitely generated} if there is a finite set $S$ with $M = \angles{S}$.
        \item We can form direct sums $ \bigoplus _{i \in I}M_i$ and \textbf{direct products} $\prod _{i \in I}M_i$ of a family $\{M_i\}_{i \in I}$ of $A$-modules.
    \end{enumerate}
\end{defn}

\begin{ex}
    The $A$-submodules of $A$ are the left ideals in $A$. In examples (2) and (3) from before, we get linear subspaces and subgroups, respectively.
\end{ex}

\begin{notebox*}[Exercise]
    let $f:A \to B$ be a ring homomorphism and $N$ a $B$-module. If $N$ is finitely generated as $B$-module, and $B$ is finitely generated as an $A$-module, then $N$ is finitely generated as an $A$ module.
\end{notebox*}

\section{Algebras}

\begin{defn}
    Let $f:A \to B$ be a ring homomorphism, and assume $A,B$ are commutative. If $a \in A$ and $b \in B$, we define a product $a \cdot b := f(a)\cdot b$. We've seen this turns $B$ into an $A$-module. But, $B$ is also a ring, and the two structures are compatible with each other in an obvious sense. The ring $B$ equipped with this $A$-module structure is said to be an $A$-algebra. So, an $A$-\textbf{algebra} is a ring $B$ together with a ring homomorphism $f:A \to B$.
\end{defn}

\begin{rmk}
    \begin{enumerate}
        \item If $A=F$ is a field and $B \neq 0$, then $f$ is injective, so $F$ can be identified with its image in $B$. So, an $F$-algebra, for $F$ a field, is a ring $B$ which contains $F$ as a subring.
        \item Any ring $A$ is automatically a $\Z$-algebra, through the unique map $\Z \to A$.
    \end{enumerate}
\end{rmk}

\section{Free Modules}

\begin{defn}
    Given a ring $A$, one can form the $A$-modules $A^k = A \oplus \overset{k }{\cdots }\oplus A$ or $\bigoplus_{i \in I}Ae_i$ for a family of elements $\{e_i\}_{i \in I}$. When $A$ is a field, this exhausts the list of all $A$-modules up to isomorphism.
\end{defn}

\begin{defn}
    Given an $A$-module $M$, we define \textbf{linearly independent} subsets of $M$ in the same way as for vector spaces. A set $S \subset M$ is a \textbf{basis} of $M$ if $S$ is linearly independent, and $M = \angles{S}$.
\end{defn}

\begin{ex}
    The \textbf{standard basis} $e_1,\dots,e_n$ is a basis of $A^n$.
\end{ex}

\begin{defn}
    An $A$-module $M$ is \textbf{free} if it has a basis $S$. In this case, we say $M$ is \textbf{free on $S$}, and there exists an isomorphism $M \cong \bigoplus_{s \in S}As,$ which is $\cong A^{|S|}$ when $S$ is finite.

    Equivalently, each element $m \in M$ has a unique expression $m = \sum a_is_i$ as a finite sum, for $a_i \in A,s_i \in S$.
\end{defn}

\begin{thm}
    We have a universal property of the pair $(M,S)$ for $S$ a basis: Given $F$ free on $S$, a map $\varphi : S \to F$, there exists a unique module homomorphism $\tilde f$ such that $f = \tilde f \circ \varphi$.
\end{thm}

% https://q.uiver.app/#q=WzAsMyxbMCwwLCJTIl0sWzIsMCwiRiJdLFsxLDIsIk0iXSxbMCwxLCJcXHZhcnBoaSIsMCx7InN0eWxlIjp7InRhaWwiOnsibmFtZSI6Imhvb2siLCJzaWRlIjoidG9wIn19fV0sWzAsMiwiZiIsMl0sWzEsMiwiXFxleGlzdHMgISBcXHRpbGRlIGYiLDAseyJzdHlsZSI6eyJib2R5Ijp7Im5hbWUiOiJkYXNoZWQifX19XV0=
\[\begin{tikzcd}
	S && F \\
	\\
	& M
	\arrow["\varphi", hook, from=1-1, to=1-3]
	\arrow["f"', from=1-1, to=3-2]
	\arrow["{\exists ! \tilde f}", dashed, from=1-3, to=3-2]
\end{tikzcd}\]

\begin{ex}
    \begin{enumerate}
        \item Any vector space $V$ over a field $F$ is a free $F$-module.
        \item $\Z^n$ is a free $\Z$-module, but $\Z/n\Z$ is not.
        \item There exists an example of a non-commutative ring $A$ with $A^n \cong A^m$ as $A$-modules, but $n \neq m$. We'll see that this never happens when $A$ is commutative.
        \item $(\Q,+)$ is not a free $\Z$-module. [Indeed, given $y \in \Q$, the equation $2x=y$ always has a solution in $x$. But, this property fails for any free $\Z$-module $\bigoplus_{i \in I}\Z e_i$]
    \end{enumerate}
\end{ex}

\begin{prop}
    Let $R$ be a non-zero commutative ring. If $M$ is a free $R$-module, with a basis of $n$ elements, then any other basis of $M$ has $n$ elements. So, we can define $\dim _R(M) := n$.
\end{prop}

\begin{proof}
    If $I$ is an ideal of $R$, we get a submodule $IM$ of $M$, generated by all products $am$, $a \in I,m \in M$. If $e_1,\dots,e_n$ is a basis of $M$, then $IM \equiv \{ (a_1e_1,\dots,a_ne_n) \mid a_i \in I\}$. Also, $M/IM$ is an $R/I$-module via the product $(r+I)(m+IM):= rm+IM$. Now, let $m$ be any maximal ideal of $R$. Then, $M / mM$ is a $k$-vector space, where $k:= R/m$. Moreover, if $\bar {e_i}=e_i+mM,1 \leq i \leq n$, then $\bar e_1,\dots,\bar e_n$ is a basis of $M/mM$ over $k$ (for linear independence, note that if $\sum r_i\bar e_i \in mM$, then $r_i \in m$ $\forall i$). The result then follows from the vector space case.
\end{proof}

\section{Exact Sequences}
\begin{defn}
    Suppose $M$ is a module, $N \subset M$ a submodule. We then have a sequence $N \overset{i}{\hookrightarrow} M \overset{\pi}{\to} M/N$ such that $\ker \pi = \Img i \equiv N$. More generally, a sequence of $A$-modules $M' \overset{\alpha}{\to}M \overset{\beta}{\to}M''$ and homomorphisms is \textbf{exact at} $M$ if $\Img \alpha = \ker \beta$. Note that $0 \to M' \overset{\alpha}{\to} M$ is exact $\iff$ $\alpha$ is injective, and $M \overset{\beta}{\to}M''\to 0$ is exact $\iff$ $\beta$ is surjective. A diagram
    $$M_1 \overset{a_1}{\to}M_2 \overset{a_2}{\to}M_3\to \cdots \to \overset{a_{n-1}}{\to} M_n$$
    is \textbf{exact} if each consecutive triplet is exact. A \textbf{short exact sequence} is an exact sequence of the form
    $$0 \to M' \overset{\alpha}{\to}M \overset{\beta}{\to}M'' \to 0,$$
    in which case $\alpha$ is injective, $\beta$ is surjective, and $M'' \cong M/M'$.
\end{defn}

\begin{ex}
    Given a finitely generated $A$-module $M$, with generators $s_1,\dots,s_r$, we have a morphism $F:= A^r = \bigoplus_{i=1}^rAs_i \overset{\varphi}{\to}M \to 0$. If $K = \ker \varphi$, we get the short exact sequence $0 \to K \to F \to M \to 0$. $K \subset F$ is the submodule of \textbf{relations} among the generators $s_i$. The short exact sequence gives a \textbf{presentation} of $M$
\end{ex}

\section{Tensor Products of Modules}

Let $A$ be a commutative ring, and $M,N$ be $A$-modules. We want to study \textbf{bilinear} maps $\varphi: M \times N \to G$, where $G$ is another $A$-module, i.e. maps $\varphi$ which are linear in each variable, when the other is fixed. The idea is to form a new module $M \otimes N$ such that the $A$-linear maps $M \otimes N \to G$ correspond to the bilinear maps $\varphi$.

\begin{ex}[First Motivation for Tensor Products]
    Suppose $V$ and $W$ are finite dimensional $F$-vector spaces. We want to study the vector space $M=\Bil (V\times W, F)$ of all bilinear maps $V \times W \to F$. Given $f \in V^*, g \in W^*$, we can form $f \otimes g \in M$ by $(f \otimes g)(x,y) := f(x)g(y)$. Note that $f \otimes g + f'\otimes g' \neq f'' \otimes g''$ in general. So the set of all maps $f \otimes g$ dose not exhaust $M$. However, we'll show
    $$M = \brax{\sum_{i=1}^n f_i \otimes g_i \mid f_i \in V^*,g_i \in W^*}=: V^* \otimes W^*$$
\end{ex}

\begin{ex}[Second Motivation for Tensor Products]
    Suppose $V$ is a finite dimensional real vector space, and $T:V \to V$ a linear operator such that $T^2+I=0$. Then, $T$ has no real eigenvalues. We may want to enlarge $V$ to get a vector space $V_\C$ over $\C$, with an extended map $T_\C :V_\C \to V_\C$ to rectify this. We do this by extension of scalars, via the inclusion $\R \hookrightarrow \C$. $V$ and $\C$ are both $\R$-modules, so we can form $V_\C := \C \otimes_\R V$, and $R$-module. But, $V_\C$ is also naturally a $\C$ module, via the operation $\lambda(z \otimes v)=(\lambda z) \otimes v$, for $\lambda,z \in \C,v \in V$. Moreover, $T$ extends to $T_\C := 1 \otimes T: V_\C \to V_\C$, which is $\C$-linear.
\end{ex}

Let $R$ be a commutative ring, $R \neq 0$. Given $M,N$ two $R$-modules, the \textbf{tensor product} $M \otimes _RN$ is an $R$-module, together with a bilinear map $\varphi :M \times N \to M\otimes N$. That is,
$$\varphi (\lambda m + m', n) = \lambda \varphi (m,n) + \varphi(m',n)\:\:\: \varphi(m, \lambda n+n')=\lambda\varphi (m,n) + \varphi(m,n').$$
We want $\varphi$ to be universal in the following sense: for any \textbf{bilinear map} $f :M \times N \to G$ of $R$-modules, there is a unique $R$-linear map $\tilde f: M \otimes_R N \to G$ such that $\tilde f \circ \varphi = f$. That is, the below diagram commutes.

% https://q.uiver.app/#q=WzAsMyxbMCwxLCJNIFxcdGltZXMgTiJdLFsyLDAsIk0gXFxvdGltZXMgX1JOIl0sWzIsMiwiRyJdLFswLDEsIlxcdmFycGhpIl0sWzAsMiwiZiIsMl0sWzEsMiwiXFxleGlzdHMgISBcXHRpbGRlIGYiLDAseyJzdHlsZSI6eyJib2R5Ijp7Im5hbWUiOiJkYXNoZWQifX19XV0=
\[\begin{tikzcd}
	&& {M \otimes _RN} \\
	{M \times N} \\
	&& G
	\arrow["{\exists ! \tilde f}", dashed, from=1-3, to=3-3]
	\arrow["\varphi", from=2-1, to=1-3]
	\arrow["f"', from=2-1, to=3-3]
\end{tikzcd}\]

The definition extends to \textbf{multilinear maps} to obtain any finite tensor product $E_1 \otimes E_2 \otimes \cdots \otimes E_n$ of $R$-modules. We construct the tensor product below.

In the following construction, we will work with $2$ $R$-modules $E$ and $F$, and construct $E \otimes _RF$. A similar idea is used to construct any finite tensor product.

\begin{defn}
    Let $E,F$ be two $R$-modules. Let $M$ be the free $R$-module generated by the \underline{set} $E \times F$. That is,
    $$M:= \bigoplus_{e \in E,f\in F}R\cdot (e,f)$$
    Let $N \subset M$ be the submodule generated by all elements of the form
    $$(\lambda e+e',f)-\lambda(e,f)-(e',f),\text{ and}$$
    $$(e,\lambda f+f') - \lambda(e,f) - (e,f')$$
    where $\lambda \in R,e,e'\in E,f,f' \in F$. We have an injection $E \times F \hookrightarrow M$ of sets. Compose this with the canonical map $M \to M/N$ to get a map $\varphi :E \times F \to M/N$. That is, $(e,f) \mapsto (e,f)+N$. We claim $\varphi$ is bilinear and a \textbf{tensor product} in the sense we want. So, we can define $E \otimes _RF$ to be $M/N$.
\end{defn}
\begin{proof}
    Bilinearity of $\varphi$ is obvious, using the relations in $N$. Now, let $f:E \times F \to G$ be any bilinear map. By the universal property of the free module generated by $E \times F$, we have an induced map $f': M \to G$ which makes the first diagram below commute. Since $f$ is bilinear, $f'$ is zero on $N$. Hence, by the universal property of quotient modules, $f'$ can be factored through $M/N$. This gives a map $\tilde f:M /N \to G$ which makes the second diagram commute. Since $M/N$ is generated by the image $\Img \varphi$ of $\varphi$, it follows that $\tilde f$ is uniquely determined. This completes the construction, and we denote $M/N$ by $E \otimes _R F$, or $E \otimes F$.
\end{proof}
% https://q.uiver.app/#q=WzAsNixbMCwxLCJFIFxcdGltZXMgRiJdLFsyLDAsIk0iXSxbMiwyLCJHIl0sWzMsMSwiRSBcXHRpbWVzIEYiXSxbNSwwLCJNL04iXSxbNSwyLCJHIl0sWzAsMiwiZiIsMl0sWzAsMV0sWzEsMiwiZiciXSxbMyw0LCJcXHZhcnBoaSJdLFszLDUsImYiLDJdLFs0LDUsIlxcdGlsZGUgZiJdXQ==
\[\begin{tikzcd}
	&& M &&& {M/N} \\
	{E \times F} &&& {E \times F} \\
	&& G &&& G
	\arrow["{f'}", from=1-3, to=3-3]
	\arrow["{\tilde f}", from=1-6, to=3-6]
	\arrow[from=2-1, to=1-3]
	\arrow["f"', from=2-1, to=3-3]
	\arrow["\varphi", from=2-4, to=1-6]
	\arrow["f"', from=2-4, to=3-6]
\end{tikzcd}\]

\begin{rmk}
    \begin{enumerate}
        \item If $x _i \in E_i$ for $1 \leq i \leq n$, write $x_1 \otimes \cdots \otimes x_n$ for the element $\varphi(x_1,\dots,x_n)$. This is called a \textbf{tensor}. We have
        $$x_1\otimes \cdots \otimes (\lambda x_i + x_i')\otimes \cdots \otimes x_n = \lambda \cdot  x_1\otimes \cdots \otimes x_n + x_1 \otimes\cdots \otimes x_i'\otimes \cdots \otimes x_n$$
        \item Every element of $E \otimes F$ is a sum of terms $x \otimes y$ with $x \in E, y \in F$, and $\lambda( x \otimes y)= (\lambda x)\otimes y = x \otimes (\lambda y)$, for any $\lambda \in R$.
    \end{enumerate}
\end{rmk}

\begin{ex}
    The tensor product of modules can involve a great deal of collapsing. For example, if $m, n \in \Z_{>0}$ with $(m,n)=1$, then $\Z/m\Z \otimes \Z/n\Z = 0$: If $x \otimes y \in \Z/m\Z \otimes \Z/n\Z$, then $n(x \otimes y) = x \otimes (ny) = x \otimes 0 = 0$, and the same for $m$. We deduce that $x \otimes y$.
\end{ex}

\begin{notebox*}[Exercise]
    Show more generally that
    $$\Z/m\Z \otimes _\Z \Z/n\Z \cong \Z/(m,n)\Z$$
\end{notebox*}

\begin{thm}
    Let $E,F,G$ be three $R$-modules. We have canonical isomorphisms
    $$\Hom (E \otimes F,G) \cong \Bil (E\times F,G) \cong \Hom (E,\Hom(F,G))$$
\end{thm}

\begin{proof}
    All the isomorphisms involved are very natural. The isomorphism $\Bil (E \times F,G) \xrightarrow{\sim}\Hom (E \otimes F,G)$ which maps $f \mapsto \tilde f$ is defined by the universal property of $E \otimes F$. It is injective by uniqueness of $\tilde f$, and surjective because composing any linear map $g : E \otimes F \to G$ with the canonical bilinear map $\varphi :E \times F \to E \otimes F$ gives a bilinear map $E \times F \to G$.

    For the second isomorphism, construct $i$ and $j$ as follows:
    $$i:\Bil (E \times F,G) \to \Hom(E,\Hom (F,G))$$
    Given $f:E \times F \to G$ and $x \in E$, we obtain a linear map $f_x: F \to G$ by $f_x(y) := f(x,y)$. Moreover, the map $x \mapsto f_x$ is linear in $x$, as desired. That is, $i$ maps $f \mapsto (x \mapsto f_x)$.
    $$j: \Hom(E,\Hom(F,G)) \to \Bil(E \times F,G)$$
    Let $\alpha:E \to \Hom(F,G)$ be $R$-linear. Define a bilinear map $f_\alpha : E \times F \to G$ by $f_\alpha(x,y) := \alpha(x)(y)$. Then, $j$ is the mapping $\alpha \mapsto f_\alpha$.

    One can check that $j$ is the inverse map to $i$.
\end{proof}

\begin{ex}
    Let $E$ and $F$ be vector spaces over the field $k$ ($\equiv R$) and set $G:= k$. Then, we have isomorphisms
    $$(E \otimes F)^* \cong \Bil (E\times F , k) \cong \Hom (E,F^*)$$
\end{ex}

\begin{cor}
    Let $E$ and $F$ be finite dimensional vector spaces over a field $k$. Then, $E \otimes F$ is finite dimensional, and $\dim_K(E \otimes F) = \dim E \cdot \dim F$. Moreover, if $e_1,\dots,e_m$ is a basis of $E$, and $f_1,\dots,f_n$ is a basis of $F$, then $\{e_i \otimes f_j\} _{1\leq i \leq m,1 \leq j \leq n}$ is a basis of $E \otimes F$.
\end{cor}

\begin{proof}
    For the last assertion, we claim that $e_i \otimes f_j$ for $1 \leq i \leq m$ and $1 \leq j \leq n$ span $E \otimes F$. Indeed, clearly $\Img \varphi = \{ v \otimes w \mid v \in V,w \in W\}$ lies in $\Span \angles{e_i \otimes f_j}$, and we know $\Img \varphi$ generated $E \otimes _kF$
\end{proof}

\begin{notebox*}[Exercise]
    With $E$ and $F$ as in the corollary, we have
    $$(E \times F)^* \cong E^* \otimes F^*$$
    Hence
    $$E \otimes F \cong \Hom(E^*,F).$$
    Under the latter isomorphism, the tensor $e \otimes f$ maps to the homomorphism $g \mapsto g(e)\cdot f$, for $g \in E^*$.
\end{notebox*}

\begin{ex}
    \begin{enumerate}
        \item (Tensor Product of 2 Homomorphisms) Let $f:M \to M'$ and $g:N \to N'$ be $R$-module homomorphisms. Then, there exists a unique $R$-linear map $f \otimes g:M \otimes N \to M ' \otimes N'$ such that $(f \otimes g)(x \otimes y) = f(x) \otimes g(y)$ for all $x \in M,y \in N$. (First, define the bilinear map $M \times N \to M' \otimes N'$ sending $(x,y) \mapsto f(x)\otimes g(y)$, and then use the universal property of $M \otimes N$)
        \item (Change of Base) Let $f : R \to S$ be any ring homomorphism, with $R$ and $S$ commutative. Recall that $S$ becomes an $R$-module by setting $r \cdot s := f(r) s$, $r \in R,s \in S$. If $M$ is any $R$-module, then $S \otimes _RM$ is an $S$-module via $s \cdot (s'\otimes m) := (ss') \otimes m$. We say that $S \otimes _RM$ is obtained from $M$ by \textbf{changing the base} from $R$ to S. A typical case of this is when $R$ is a subring of $S$, in which case $S \otimes _R M$ is an \textbf{extension of scalars} from $R$ to $S$. For example, if $R \subset S$ is a subring, then $S \otimes _R R[x_1,\dots,x_n]\cong S[x_1,\dots,x_n]$.
    \end{enumerate}
\end{ex}

\begin{prop}
    $$E \otimes _R F \cong F \otimes _R E$$
    as modules.
\end{prop}

\begin{proof}
    The map $E \times F \to F \otimes _RE$ which maps $(x,y) \mapsto y \otimes x$ is bilinear, so by the universal property there exists an $R$-linear map
    $$E \otimes _RF \to F \otimes _RE$$
    sending $x \otimes y \mapsto y \otimes x$, which has a clear inverse by symmetry.
\end{proof}

\begin{prop}
    The (unique) map $(x \otimes y) \otimes z \mapsto x \otimes (y \otimes z)$ induces an isomorphism
    $$(E \otimes F) \otimes G \cong E \otimes (F \otimes G)$$
\end{prop}

\begin{proof}
    Since the tensors $(x \otimes y) \otimes z$ generate $(E \otimes F) \otimes G$, the uniqueness is clear. To prove existence, fix a $z \in G$, then the map $\lambda_z: E \times F \to E \otimes (F \otimes G)$ mapping
    $$\lambda_z(x,y) := x \otimes (y \otimes z)$$
    is bilinear in $x$ and $y$, so induces a linear map $\tilde \lambda _z : E \otimes F \to E \otimes (F \otimes G)$. Furthermore, consider the map $(E \otimes F) \times G \to E \otimes (F \otimes G)$ which sends $(\alpha, z) \mapsto \tilde \lambda _z(\alpha)$, and note that it's bilinear, hence factors through to get the desired linear isomorphism.
\end{proof}

\begin{ex}
    \begin{enumerate}
        \item Let $F$ be a field and $V$ and $W$ be $F$-vector spaces with $\dim_F V=n,\dim _FW = m$. Then, if $E := (e_1,\dots,e_n)$ and $E' := (e_1',\dots,e_m')$ are ordered bases of $V,W$ respectively, then
        $$\{ e_i \otimes e_j' \mid 1 \leq i \leq n, 1 \leq j \leq m\}$$
        forms a basis of $V \otimes W$. Now, suppose $S:V \to V$ and $T:W \to W$ are $F$-linear maps represented by the matrices $A=(S)_E,B = (T)_{E'}$ respectively. We've seen that $S \otimes T : V \otimes W \to V \otimes W$ is a linear map characterized by
        $$(S \otimes T)(v \otimes w)= S(v) \otimes T(w).$$
        If the basis $\{ e_i \otimes e_j'\}$ is ordered as follows:
        $$e_1\otimes e_1',\dots,e_1 \otimes e_m',\dots,e_n \otimes e_1',\dots,e_n\otimes e_m',$$
        then the matrix representing $S \otimes T$ is denoted by $A \otimes B$, and has the block form
        $$\matr{a _{1,1}B & \cdots & a_{1,n}B \\ \vdots & \ddots & \vdots \\ a_{n,1}B & \cdots & a_{n,n}B}=\{ a_{ij}B\}_{1\leq i,j \leq n}.$$
        This is called the \textbf{Kronecker product} or \textbf{tensor product} of the matrices $A$ and $B$.
        \item \textbf{Tensor product of $R$-algebras:} Let $R$ be a commutative ring, and $A$ and $B$ be two commutative $R$-algebras, via the ring homomorphisms $f:R \to A,g : R \to B$. Since $A$ and $B$ are $R$-modules, we can form their tensor product $ D:=A \otimes _R B$, which is an $R$-module. We now define a product on $D$: Consider the map $A \times B \times A \times B \to D$ mapping $(a,b,a',b,') \mapsto aa' \otimes bb'$. This is $R$-linear in each factor, hence determines an $R$-module homomorphism
        $$A \otimes B \otimes A \otimes B \to D,$$
        and hence an $R$-module homomorphism $D \otimes D \to D$. This in turn corresponds to a bilinear mapping $\mu : D \times D \to D$ such that $\mu (a \otimes b, a' \otimes b') = (aa')\otimes (bb')$. Of course, we could've written this formula directly, but this construction shows $\mu$ is well defined. We have thus defined a product on $D = A \otimes _R B$. On elements of the form $a \otimes b$, it looks like $(a \otimes b)(a' \otimes b') = (a a') \otimes (bb')$, and in general
        $$\paren{\sum_i a_i \otimes b_i}\paren{\sum_j a_j' \otimes b_j'}= \sum_{i,j} a_ia_j' \otimes b_ib_j'.$$
        This makes $D$ into a commutative ring with identity $1 \otimes 1$. Also, $D$ is an $R$-algebra via the map $r \mapsto f(r) \otimes g(r)$ (a ring homomorphism $R \to D)$. These constructions  also extend to the non-commutative setting.
    \end{enumerate}
\end{ex}

\begin{defn}[Tensor Products as $R$-Algebras]
    For $R$ commutative, we say that an \textbf{$R$-algebra} is a (possibly noncommutative) ring $A$ together with a ring homomorphism $f:R \to A$ such that the subring $\Img (f)$ of $A$ is contained in the center $Z(A)$ of $A$. Equivalently, $A$ is a ring that is also a left $R$-module such that $r(ab) = (ra)b=a(rb)$ for all $r \in R, ab \in A$ (In this case, $f:R \to A$ sends $r$ to $r \cdot 1_A$, and one can check that $f(R)$ is in the center). Now, if $A$ and $B$ are $R$-algebras, then the product $(a \otimes a')(a' \otimes b') := (aa') \otimes (bb')$ is well defined and makes $A \otimes _R B$ into an $R$-algebra.
\end{defn}

\begin{prop}
    Suppose $R$ is a commutative ring, and $I,J$ are ideals of $R$. Then,
    $$(R/I) \otimes _R (R/J) \cong R/(I+J)$$
    as $R$-algebras.
\end{prop}

\begin{proof}
    Define $f : R/I \times R/J \to R/(I+J)$ by $f(x+I,y+J):= xy +I+J$. $f$ is $R$-bilinear. Moreover, for any $R$-module $G$ and $R$-bilinear $g : R/I \times R/J \to G$, define $\tilde g : R/(I+J) \to G$ by $\tilde g(x+I+J) = g(x+I,1+J)$. Then,
    \begin{enumerate}
        \item $\tilde g$ is well defined: If $x+I+J = x'+I+J$, then $x'=x+a+b$ for $a\in I, b \in J$, hence $x'+I=x+b+I$, and
        $$g(x'+I,1+J)=g(x+b+I,1+J)=g(x+I,1+J)+g(b+I,1+J)$$
        $$=g(x+I,1+J)+g(1+I,b+J)=g(x+I,1+J)$$
        \item $\tilde g$ is $R$-linear: clear
        \item $\tilde g$ satisfies the universal property:
        $$\tilde g \circ f(x+I,y+J) = \tilde g(xy+I+J)=g(xy+I,1+J)=g(x+I,y+J),$$
        hence $\tilde g$ is uniquely determined, since $f$ maps $R/I \times R/J$ onto $R/I+J$.
    \end{enumerate}
    Therefore, $R/I+J$ and $R/I \otimes _R R/J$ are isomorphic as $R$-modules. Moreover, the map $x+I+J \mapsto (x+I) \otimes (1+J)$ is an explicit isomorphism between them, and is an $R$-algebra isomorphism.
\end{proof}

\begin{cor}
    $$\Z/m\Z \otimes _\Z \Z / n\Z \cong \Z/(m,n)\Z$$
\end{cor}

\begin{prop}
    Let $N:= \bigoplus _{i=1}^nN_i$ be a direct sum of $R$-modules. Then, we have an isomorphism
    $$M \otimes N \cong \bigoplus_{i=1}^n (M \otimes N_i)$$
    as $R$-modules.
\end{prop}

\begin{proof}
    Consider the map $M \times (\bigoplus_{i=1}^n N_i) \to \bigoplus_{i=1}^n(M \otimes N_i)$ mapping
    $$(m, \sum n_i) \mapsto \sum m \otimes n_i.$$
    It is $R$-bilinear, hence induces a map
    $$f : M \otimes (\bigoplus_{i=1}^nN_i) \to \bigoplus_{i=1}^n (M \otimes N_i).$$
    In the other direction, for each $i \in [1,n]$, we have a map $M \times N_i \to M \otimes (\bigoplus_{i=1}^nN_i)$ sending $(m,n_i) \to m \otimes n_i$. These induce maps
    $$M \otimes N_i \to M \otimes (\bigoplus_{i=1}^nN_i)$$
    for each $i$. By the universal property of the direct sum, we obtain a map
    $$g: \bigoplus_{i=1}^n (M \otimes N_i) \to M \otimes (\bigoplus_{i=1}^n N_i),$$
    and one can check that $f$ and $g$ are inverses of each other.
\end{proof}

\begin{notebox*}[Exercise]
    Show that the previous proposition holds for arbitrary direct sums:
    $$M \otimes (\bigoplus_{i \in I}N_i) \cong \bigoplus_{i \in I}(M \otimes N_i).$$
    We say that $\otimes$ \textbf{commutes with direct sums.}
\end{notebox*}

The next result shows that $E$ is a free $R$-module, there is no collapsing in a tensor product $F \otimes _RE$.

\begin{prop}
    Let $E$ be free over $R$ with basis $\{ e_i\}_{i \in I}$. Then, for any $R$-module $F$, each element of $F \otimes E$ has a unique expression of the form
    $$x = \sum_{i \in I}f _i \otimes e_i,$$
    where only finitely many $f_i \in F$ are non-zero.
\end{prop}

\begin{proof}
    Assume first that $\dim _R E=1$, so $E = Re$ for some $e \in E$. Since $f \otimes (\lambda e)= (\lambda f) \otimes e$, and any sum
    $$\sum_{i=1}^n (f_i \otimes e)= \paren{\sum_{i=1}^n f_i}\otimes e$$
    by linearity, we deduce that any element of $F \otimes E$ may be written as $f \otimes e$ for some $f \in F$. The linear map $F \to F \otimes E$ sending $f \mapsto f \otimes e$ is an isomorphism with inverse $F \otimes E \to F$ mapping $f \otimes (\lambda e) \to \lambda f$, for $f \in F$, $\lambda \in R$. Since $F \otimes E \cong F$, each element of $F \otimes E$ has a unique expression as $f \otimes e$ for some $f \in F$.

    Now, in general, if $E$ is free with basis $\{e_i\}_{i \in I}$, then we have
    $$F \otimes E = F \otimes (\bigoplus_{i \in I}Re_i) = \bigoplus_{i \in I}(F \otimes Re_i),$$
    and the result follows from the rank $1$ case.
\end{proof}

\begin{cor}
    Let $E,F$ be free $R$-modules with bases $\{e_i\}_{i \in I}$ and $\{f_j\} _{j \in J}$ respectively. Then, $E \otimes F$ is free with basis $\{e_i \otimes f_j\}_{i\in I,j \in J}$. If $E$ and $F$ are both finite dimensional, then we have
    $$\dim _R(E \otimes F) = \dim _R(E) \cdot \dim_R(F)$$
\end{cor}

\begin{rmk}
    When $E$ is free over $R$, every element of $F \otimes E$ can be viewed as a "formal" linear combination of elements in a basis $\{e_i\}$ of $E$, with coefficients in $F$. In particular
    $$R \otimes E \cong E \otimes R \cong E$$
    via the map $r \otimes e \mapsto re$. (Check: Is $E$ being free here necessary?)
\end{rmk}

\begin{ex}
    If $V$ is a vector space over a field $F$, and $K \supset F$ is a field extension, then $V_K := K \otimes _F V$ is a $K$-vector space. If $\dim _F (V) =n$, and $\{x_1,\dots,x_n\}$ is an $F$-basis of $V$, then observe
    $$V_K = K \otimes _F V \cong K \otimes _F(F^n) \cong (K \otimes _F F)^n \cong K^n$$
    and $\{ 1\otimes x_1,\dots, 1 \otimes x_n\}$ is a $K$-basis for $V_K$.

    Explicitly, given
    $$x = \sum_{j=1}^m a_j \otimes v_j \in V_K,$$
    we may write each $v_j = \sum_{i=1}^n\lambda_{i,j}x_i$ for $\lambda_{i,j} \in F$, and therefore
    $$x = \sum_{j=1}^m a_j \otimes \paren{\sum_{i=1}^n\lambda_{i,j}x_i}= \sum_{i=1}^n \paren{\sum_{j=1}^m \lambda _{i,j} a_j}(1 \otimes x_i).$$
\end{ex}

\begin{notebox*}[Exercise]
    Show that if $K \supset F$ is a finite extension, and $[K:F] =\dim _FK$, then
    $$\dim _FV_k = [K:F]\dim _FV.$$
\end{notebox*}

\begin{cor}
    Let $f:R \to S$ be a ring homomorphism. Then,
    $$S \otimes _R R^n \cong S^n$$
    as $S$-modules. This is called \textbf{extension of scalars for free modules.}
\end{cor}

\begin{proof}
    Use the isomorphism $S\otimes _R R \cong S$.
\end{proof}

\section{Graded Rings and Algebras}

\begin{defn}
    A ring $A$ is \textbf{graded} if $A$ is a direct sum
    $$A = \bigoplus_{i = 0}^\infty A_i$$
    of additive subgroups $A_i$ such that the ring multiplication maps $A_i \times A_j \hookrightarrow A_{i+j}$ for all $i,j\geq0$. In particular, we see that $A_0$ is a subring of $A$. The elements of $A_k$ are called the \textbf{homogeneous elements of degree $k$}.

    A \textbf{graded ring homomorphism} $f:A \to A'$ is a homomorphism of rings that respects the grading, i.e. $f(A_k) \subset A'_k$ for all $k$.

    The definition of \textbf{graded $R$-algebras} is similar - in this case each $A_k$ is an $R$-submodule of $A$, where $R$ is the ring. (Usually we have $R \equiv A_0$)
\end{defn}

\begin{ex}
    For $R$ a commutative ring,
    $$R[x_1,\dots,x_n] = \bigoplus _{k=0}^\infty R[x_1,\dots,x_n]_k,$$
    where $R[x_1,\dots,x_n]_k :=$ homogeneous polynomials of degree $k$.
\end{ex}

\section{The Tensor Algebra}

\begin{defn}
    Let $R$ be a commutative ring, and $M$ be an $R$-module. We then define $T^0M:=R,T^1M:=M$, and
    $$T^kM := M \otimes _R \overset{k}{\cdots} \otimes _R M $$
    The elements of $T^kM$ are called $k$-tensors.
\end{defn}

\begin{defn}
    The \textbf{tensor algebra} $TM$ is the graded $R$-algebra
    $$ TM = \bigoplus_{k=0}^\infty T^kM$$
    as an $R$-module. The product in $TM$ is induced by the maps $T^m M \otimes T^nM \to T^{m+n}M$ sending $(x_1\otimes \cdots \otimes x_m) \otimes (y_1\otimes \cdots \otimes y_n)$ to $x_1\otimes \cdots \otimes x_m \otimes y_1 \otimes \cdots \otimes y_m$ for $m,n \geq 0$, which extend to a map $TM \otimes TM \to TM$ using the universal properties.

    The identification of $M$ with $T^1M$ gives an inclusion $i : M \hookrightarrow TM$. The pair $(TM,i)$ is characterized by the following universal property:

    For any $R$-algebra $A$ and $R$-linear map $j:M \to A$ (of $R$-modules), there exists a unique map $\tilde j : TM \to A$ of $R$-algebras such that $\tilde j \circ i = j$
\end{defn}
% https://q.uiver.app/#q=WzAsMyxbMCwxLCJNIl0sWzIsMCwiVE0iXSxbMiwyLCJBIl0sWzAsMiwiaiIsMl0sWzAsMSwiaSIsMCx7InN0eWxlIjp7InRhaWwiOnsibmFtZSI6Imhvb2siLCJzaWRlIjoidG9wIn19fV0sWzEsMiwiXFxleGlzdHMgISBcXHRpbGRlIGoiLDAseyJzdHlsZSI6eyJib2R5Ijp7Im5hbWUiOiJkYXNoZWQifX19XV0=
\[\begin{tikzcd}
	&& TM \\
	M \\
	&& A
	\arrow["{\exists ! \tilde j}", dashed, from=1-3, to=3-3]
	\arrow["i", hook, from=2-1, to=1-3]
	\arrow["j"', from=2-1, to=3-3]
\end{tikzcd}\]

\begin{defn}
    Let $P$ be an algebra over $R$. We say that $P$ is a \textbf{non-commutative polynomial algebra} if there exist elements $x_1,\dots,x_n \in P$ such that the elements $x_I:= x_{i_1}\cdots x_{i_s},I = (i_1,\dots,i_s), 1 \leq i_k \leq n$ form a basis of $P$ over $R$. We call these elements $x_I$ \textbf{non-commutative monomials} in $x_1,\dots,x_n$. By convention, when $s=0$, the corresponding monomial is the identity element in $P$.

    We say that $x_1,\dots,x_n$ are independent noncommutative variables over $R$, and write $P \cong R\angles {x_1,\dots,x_n}$, the latter being the non-commutative polynomial ring in these variables.

    In other words, a non-commutative polynomial algebra in $x_1,\dots,x_n$ is almost the same as $R[x_1,\dots,x_n]$, except that the variables don't commute.
\end{defn}

\begin{prop}
    Let $E$ be a free $R$-module, $\dim _RE =n$. Then, $TE$ is isomorphic to the non-commutative polynomial algebra on $n$ variables over $R$. In other words, if $e_1,\dots,e_n$ is a basis of $E$ over $R$, then the elements
    $$e_I := e_{i_1}\otimes \cdots \otimes e_{i_s},I=(i_1,\dots,i_s),1\leq i_k \leq n$$
    form a basis of $T^sE$, and each element of $TE$ has a unique expression as a finite sum $\sum _I r_Ie_I$ where $r_I \in R$, with all but finitely many $r_I$ equal to $0$.
\end{prop}

\begin{proof}
    Recall that for any $R$-module $F$, every element of $F \otimes E$ has a unique expression of the form $\sum_{i=1}^n f_i \otimes e_i$, with $f_i \in F$. The proposition follows immediately from this.
\end{proof}

\begin{rmk}
    The tensor algebra construction is \textbf{functorial}, in the sense that a morphism $f:M \to N$ between $R$-modules induces a map $Tf :TM \to TN$ between tensor algebras.
\end{rmk}

Now, we don't really like non-commutative polynomial rings. We've shown how a free $R$ module can be translated into a non-commutative polynomial algebra. The key problem is that $x \otimes y = y \otimes x$. Let's fix it!

\section{Symmetric Powers and Symmetric Algebras}

\begin{defn}
    Let $M$ be an $R$-module and $n \geq 0$. The \textbf{symmetric power} $S^nM \equiv \w{Sym} ^nM$ is an $R$-module together with an $R$-multilinear map $\psi :M \times \overset{n}{\cdots } \times M \to S^nM$, with $\psi (x_1,\dots,x_n)$ denoted $x_1\cdot x_2 \cdots x_n$, that is symmetric. This means that $x_{\sigma(1)}\cdots x_{\sigma (n)}=x_1\cdots x_n$ for all $\sigma \in S_n$. The pair $(S^nM, \psi)$ satisfies the universal property:

    For any $R$-module $N$, and symmetric multilinear map $f:M \times \overset{n}{\cdots} \times M \to N$, there exists a unique map $\tilde f : S^nM \to N$ of $R$-modules such that $\tilde f \circ \psi = f$.
\end{defn}
% https://q.uiver.app/#q=WzAsMyxbMCwxLCJNIFxcdGltZXMgXFxvdmVyc2V0e259e1xcY2RvdHMgfSBcXHRpbWVzIE0iXSxbMiwwLCJOIl0sWzIsMiwiU15uTSJdLFswLDIsIlxccHNpIiwyXSxbMCwxLCJmIl0sWzEsMiwiXFxleGlzdHMgISBcXHRpbGRlIGYiLDAseyJzdHlsZSI6eyJib2R5Ijp7Im5hbWUiOiJkYXNoZWQifX19XV0=
\[\begin{tikzcd}
	&& S^nM \\
	{M \times \overset{n}{\cdots } \times M} \\
	&& {N}
	\arrow["{\exists ! \tilde f}", dashed, from=1-3, to=3-3]
	\arrow["f", from=2-1, to=1-3]
	\arrow["\psi"', from=2-1, to=3-3]
\end{tikzcd}\]
\begin{proof}
    We see that $S^nM := T^nM / J$, where $J$ is the $R$-submodule of $T^nM$ generated by all $x_{\sigma(1)}\otimes \cdots \otimes x_{\sigma (n)}- x_1\otimes \cdots \otimes x_n$. (Exercise: Verify)
\end{proof}

\begin{notebox*}[Convention]
    $S^0M = R,S^1M=M$.
\end{notebox*}

We have maps $S^aM \otimes S^bM \to S^{a+b}M$, $a,b,\geq 0$ sending
$$(x_1\cdots x_a) \otimes (y_1\cdots  y_b) \mapsto x_1 \cdots x_a y_1 \cdots y_b.$$
These make $S^\cdot M= \bigoplus_{n=0}^\infty S^nM$ into a graded \textit{commutative} $R$-algebra, which is universal for $R$-module homomorphisms from $M$ to any commutative $R$-algebra $A$, in the same way as for the tensor algebra. In fact, one can define $S^\cdot M$ as a commutative quotient of $T^\cdot M$: $S^\cdot M$ is $T^\cdot M$ modded out by the ideal generated by all two sided ideals generated by all differences $x \otimes y - y \otimes x \in T^2M$, for all $x,y \in M$.

\begin{prop}
    $$S^n(M \oplus N) \cong \bigoplus_{a+b=n}S^aM \otimes _RS^bN$$
\end{prop}

\begin{proof}
    There is a map $S^aM \otimes S^bN$ to $S^{a+b}(M \oplus N)$ using the inclusions $M,N \subset M \bigoplus N$, sending $x_1\cdots x_a \otimes y_1\cdots y_b$ to $x_1\cdots x_a y_1\cdots y_b$, as before.

    In the other direction, consider the symmetric $R$-linear map
    $$f: (M \oplus N) \times \overset{n}{\cdots} \times (M \oplus N) \to \bigoplus_{a+b=n} (S^aM \otimes _R S^bN)$$
    defined by
    $$f((x_1,y_1)\times \cdots \times (x_n,y_n)) := \bigoplus_{a+b=n}\paren{\sum_{\{ i_1,\dots, i_a \}\amalg \{ j_1,\dots,j_b\} =[1,n] } x_{i_1} \cdots x_{i_a} \otimes y_{j_1}\cdots y_{j_b}}$$
    This induces the inverse map
    $$\tilde f : S^n(M \oplus N) \to \bigoplus_{a+b = n}(S^aM \otimes S^bN)$$
\end{proof}

\begin{ex}
    \begin{enumerate}
        \item Suppose that $M =R$. Then, $S^nR = T^nR \cong R$ for all $n \geq 0$
        \item (Corollary) Suppose $E$ is free over $R$ with basis $e_1,\dots,e_k$. Then $S^nE$ is free on $e_{i_1}\cdots e_{i_n}$ for all $1 \leq i_1 \leq \cdots \leq i_n \leq k$, or all monomials $e_1^{a_1}\cdots e_k^{a_k}$, $a_i \geq 0, \sum_{i}a_i =n$. In other words,
        $$S^\cdot M \cong R [e_1,\dots,e_k]$$
        is a polynomial ring.
    \end{enumerate}
\end{ex}

\begin{notebox*}[Exercise]
    Show that the isomorphism
    $$S^n (M \bigoplus N) \cong \bigoplus_{a+b=n}S^aM \otimes _R S^bN$$
    is nothing but the $n$-part of a graded isomorphism $S^\cdot (M \oplus N) \cong S^\cdot M \otimes _R S^\cdot N$ of graded $R$-algebras.
\end{notebox*}

\section{Motivation for the Symmetric Algebra}

Given a space $X = \C^n$, we have the ring $R:= \C [x_1,\dots,x_n]$ of polynomial functions on $X$. We can recover $X$ from the ring $R$ by looking at $\Specm (R):= \{$all maximal ideals in $R\}$, since any maximal ideal in $R$ is of the form $m_a := (x_1-a_1,\dots,x_n-a_n)$ for $a = (a_1,\dots,a_n) \in \C^n$.

Now, suppose $V$ is a complex $n$-dimensional vector space. How do we define $R_V$, the ring of all polynomial functions on $V$? One way: Choose a basis $e_1,\dots ,e_n$ of $V$, use this to get an isomorphism $V \cong \C^n$ of vector spaces, and use $R = \C[x_1,\dots,x_n]$ to get a definition for $V$. This isn't as satisfying, since we had to choose a basis, i.e. it's not an intrinsic definition of $R_V$.

To get a basis-free definition, note that $R$ is a graded $\C$-algebra:
$$R = \C[ x_1,\dots,x_n] = \bigoplus_{k=0}^\infty \C[x_1,\dots,x_n]_k$$
and $\C[x_1,\dots,x_n]_1$ is essentially the set of all linear forms $a_1x_1+\cdots + a_nx_n$ for $a_i \in \C$. Therefore, $\C[x_1,\dots,x_n]_1 \equiv (\C^n)^*$. But, the latter has a basis free definition for $V$, namely dual space $V^* := \Hom _\C (V,\C)$.

\begin{defn}
    For an $n$-dimensional vector space $V$,
    $$R_V := \w{Sym} (V^*)$$
    from which $V$ can be recovered as $V = \Specm (S^\cdot V^*)$
\end{defn}

This construction is useful when one studies vector bundles on manifolds, (such as the tangent bundle $TX \to X$ of a manifold $X$ where intrinsically defined objects make sense).
